% Options for packages loaded elsewhere
\PassOptionsToPackage{unicode}{hyperref}
\PassOptionsToPackage{hyphens}{url}
%
\documentclass[
  11pt,
  a4paper,
]{article}
\usepackage{amsmath,amssymb}
\usepackage{iftex}
\ifPDFTeX
  \usepackage[T1]{fontenc}
  \usepackage[utf8]{inputenc}
  \usepackage{textcomp} % provide euro and other symbols
\else % if luatex or xetex
  \usepackage{unicode-math} % this also loads fontspec
  \defaultfontfeatures{Scale=MatchLowercase}
  \defaultfontfeatures[\rmfamily]{Ligatures=TeX,Scale=1}
\fi
\usepackage{lmodern}
\ifPDFTeX\else
  % xetex/luatex font selection
\fi
% Use upquote if available, for straight quotes in verbatim environments
\IfFileExists{upquote.sty}{\usepackage{upquote}}{}
\IfFileExists{microtype.sty}{% use microtype if available
  \usepackage[]{microtype}
  \UseMicrotypeSet[protrusion]{basicmath} % disable protrusion for tt fonts
}{}
\makeatletter
\@ifundefined{KOMAClassName}{% if non-KOMA class
  \IfFileExists{parskip.sty}{%
    \usepackage{parskip}
  }{% else
    \setlength{\parindent}{0pt}
    \setlength{\parskip}{6pt plus 2pt minus 1pt}}
}{% if KOMA class
  \KOMAoptions{parskip=half}}
\makeatother
\usepackage{xcolor}
\usepackage[margin=26mm]{geometry}
\usepackage{longtable,booktabs,array}
\usepackage{calc} % for calculating minipage widths
% Correct order of tables after \paragraph or \subparagraph
\usepackage{etoolbox}
\makeatletter
\patchcmd\longtable{\par}{\if@noskipsec\mbox{}\fi\par}{}{}
\makeatother
% Allow footnotes in longtable head/foot
\IfFileExists{footnotehyper.sty}{\usepackage{footnotehyper}}{\usepackage{footnote}}
\makesavenoteenv{longtable}
\setlength{\emergencystretch}{3em} % prevent overfull lines
\providecommand{\tightlist}{%
  \setlength{\itemsep}{0pt}\setlength{\parskip}{0pt}}
\setcounter{secnumdepth}{-\maxdimen} % remove section numbering
\usepackage{lmodern,amsmath,amssymb,booktabs,longtable,array,microtype,fancyhdr,needspace}
\pagestyle{fancy}\fancyhf{}\fancyhead[L]{Residual charging and spare sources}\fancyhead[R]{Reviewer draft v9}\fancyfoot[C]{\thepage}
\setlength{\headheight}{14pt}
\setlength{\parskip}{4pt plus 1pt}
\setlength{\emergencystretch}{2em}
\AtBeginDocument{\hypersetup{pdftitle={Residual charging stability and spare sources},pdfauthor={Paul Lenz},pdfsubject={Candidate theorems; independent review open}}}
\ifLuaTeX
  \usepackage{selnolig}  % disable illegal ligatures
\fi
\usepackage{bookmark}
\IfFileExists{xurl.sty}{\usepackage{xurl}}{} % add URL line breaks if available
\urlstyle{same}
\hypersetup{
  hidelinks,
  pdfcreator={LaTeX via pandoc}}

\author{}
\date{}

\begin{document}

\section{Residual charging stability and spare sources in diameter-two
edge-critical
graphs}\label{residual-charging-stability-and-spare-sources-in-diameter-two-edge-critical-graphs}

\textbf{Paul Lenz --- research direction. Mathematical development,
software, drafting and internal checking: ChatGPT (Geeps).} 7 September
2026. Reviewer draft v9.

\textbf{Status.} Candidate mathematical theorems with explicit proofs
and internal regression checks. Independent mathematical review and
novelty assessment are OPEN. Five local logical implications have been
checked in Lean; the graph-to-data construction, cardinality arguments
and theorems below are not yet formally verified. This work changes
neither the frozen finite-order papers nor their theorem ledger.

\subsection{Abstract}\label{abstract}

We isolate a residual charging argument for diameter-two edge-critical
graphs and quantify its loss using the common source--supplement pair
budget. For \(a=n-1-\Delta\) and \(t=m-\Delta(n-\Delta)\), the inherited
charging argument gives \(t\le(3/2-\sqrt2)a^2\). An exact deficit
identity forces concentration of the label demands near its scalar
optimum. A joint pair-capacity inequality then excludes that
concentration at every \(a\ge25\), giving \(t<(3/2-\sqrt2-1/5000)a^2\).
With 24 elementary rounding cases this yields a candidate maximum-degree
threshold \(\beta_1=0.61316824745\ldots\) for all \(n\ge4\), improving
the project's previous \(0.61327045983\ldots\) coefficient. A separate
weighted theorem bounds the number of maximum-demand labels supported by
\(h+k\) maximum-residual sources, for \(0\le k\le h^2\). In particular,
two spare sources permit at most \(2h+2\) such labels for \(h\ge2\).
Neither theorem uses positive-surplus residual activity, finite-order
certificate enumerations, Fan's density estimate, or a weak-core
reduction. Numerical tests support, but do not replace, the proofs.

\subsection{1. Scope, main statements and prior
work}\label{scope-main-statements-and-prior-work}

A graph is simple and finite. Disconnected pairs have infinite distance.
A graph \(G\) is diameter-two edge-critical if its diameter is two and
deleting any edge increases the diameter. Let \(n=|V(G)|\), \(m=e(G)\)
and \(b=\Delta(G)\). The Murty--Simon conjecture asks whether
\(m\le\lfloor n^2/4\rfloor\), with equality only for the balanced
complete bipartite graph.

\textbf{Theorem A (quantitative charging loss).} In the
minimum-complement-degree construction below, let \(a=n-1-b\ge25\) and
\(t=m-b(n-b)\). Put

\[c=\frac32-\sqrt2,\qquad c_1=c-\frac1{5000}.\]

Then

\[\boxed{t<c_1a^2.} \tag{1}\]

\textbf{Corollary A (maximum degree).} For \(n\ge4\), define

\[\beta_1=\frac{1/2+\sqrt{c_1}}{1+\sqrt{c_1}}=0.6131682474535585\ldots.\]

Then

\[\boxed{\Delta(G)\ge\beta_1n\quad\Longrightarrow\quad e(G)<\lfloor n^2/4\rfloor.} \tag{2}\]

The decimal is explanatory; the radical is the exact coefficient. The
improvement from the earlier project coefficient
\(\beta_0=(10-\sqrt2)/14=0.613270459830\ldots\) is small numerically.
The new ingredient is a loss proportional to \(a^2\), rather than
strictness only at the exact scalar optimum.

\textbf{Theorem B (spare sources).} Let \(h\ge1\) be the largest
residual degree on \(B\), and let exactly \(z\) sources have residual
degree \(h\). If \(h\le z\le h^2+h\), and \(\ell\) labels have minimum
demand \(h\), then

\[\boxed{h\ell\le(2h-1)z.} \tag{3}\]

Consequently, with \(z=h+2\) and \(h\ge2\), one has \(\ell\le2h+2\).
With \(z=h+1\) one recovers \(\ell\le2h\). No positive-surplus
assumption is made. The bound is attained by reduced incidence systems
described below; sharpness for actual critical graphs is not claimed.

\textbf{Literature positioning.} Haynes, Henning, van der Merwe and Yeo
{[}1{]} prove the conjecture for \(\Delta\ge0.7n\), and for \(n\ge2000\)
with \(\Delta\ge0.6789n\). The canonical abstract of the 2016 preprint
{[}2{]} advertises \(0.676n\) for every \(n\). An author-posted full
text associated with that record states \(0.6756n\) in Theorem 3.5, with
other metadata differences detailed in
\texttt{literature/COMPARISON.md}. The earlier project comparison
mentioning only {[}1{]} was incomplete. Complement quasi-edges and
supplements are established tools, not introduced here. The present
argument charges residual degrees against selected label demands, then
uses simultaneous pair capacity; {[}2{]} uses an inductive quasi-clique
estimate. We do not claim a completed audit of {[}1{]} or {[}2{]}, a
best-known threshold, priority, or novelty. Those results are
comparisons, not proof inputs to (1)--(3).

The inherited charging argument and one-spare-source lemma are preserved
in project checkpoints {[}4--5{]}. We reproduce their needed ingredients
below, eliminating any dependency on classification of
total-domination-critical complements. The full n=25, n=27 and n=28
numerical proofs are not dependencies.

\subsection{2. A direct quasi-edge
construction}\label{a-direct-quasi-edge-construction}

Put \(H=\overline G\). An adjacent pair totally dominates a graph when
every vertex has a neighbour in the pair. For any graph \(J\), its
complement has a pair at distance greater than two exactly when \(J\)
has an adjacent total dominating pair: the two endpoints are nonadjacent
in the complement and have no common neighbour there. This remains valid
for disconnected pairs. Thus \(H\) has no such pair, while \(H+uw\) has
one whenever \(uw\) is a nonedge of \(H\).

\textbf{Lemma 2.1 (unique exception).} Suppose \(uw\) is a nonedge of
\(H\) and both endpoints miss a third vertex. There is an existing edge
\(ui\), or symmetrically \(wi\), whose open-neighbourhood union is
\(V(H)\setminus\{w\}\) in the former orientation. Write \(ui\to w\).

\emph{Proof.} An adjacent total dominating pair in \(H+uw\) cannot be
\(\{u,w\}\), because that pair still misses the third vertex. A pair
disjoint from \(\{u,w\}\) has unchanged domination, also impossible. Its
edge therefore already existed and it uses exactly one endpoint of
\(uw\). The only newly dominated vertex is the other endpoint. It is the
unique exception of that existing edge in \(H\). In particular \(u\) and
\(i\) both miss \(w\). \(\square\)

Choose \(v\) of minimum degree in \(H\), and set

\[A=N_H(v),\quad B=V(H)\setminus N_H[v],\quad a=|A|=n-1-b,\quad |B|=b.\]

Put \(C=H[A]\) and \(F=\overline C\) on \(A\). Write \(d_i=d_F(i)\).
Each missing unordered pair in \(H[B]\) misses \(v\); its quasi-edge
auxiliary lies in \(A\), to dominate \(v\). Choose exactly one such
quasi-edge for each missing pair. Call the chosen cross-edges
\emph{selected}, and all other \(H[A,B]\) edges \emph{residual}. Orient
the missing pair from its selected source \(u\) to its supplement \(w\).

Distinct missing pairs select distinct cross-edges: the \(B\)-endpoint
and unique exception recover the pair. At one source the labels and
supplements are separately distinct. The selected arcs on \(B\) form an
oriented simple graph: an unordered pair supports at most one arc, not
one in each direction.

Let \(\rho_u\) and \(R_i\) be residual degrees at \(u\in B\) and
\(i\in A\), with common sum \(r\). Let \(q_u,p_u\) be selected outdegree
and indegree, \(x_i\) the actual selected degree of label \(i\), and
\(Q=\sum q_u=\sum x_i\). Define

\[t=m-b(n-b),\qquad s_i=\max(0,d_i-R_i),\qquad S=\sum_i s_i.\]

\textbf{Ledger.} Selected cross-edges and existing \(H[B]\) edges
together number \(\binom b2\). Counting the remaining edges gives

\[e(F)=r+t,\quad \sum_i d_i=2(r+t),\quad \sum_iR_i=\sum_u\rho_u=r. \tag{4}\]

Minimum degree in \(H\) gives \(d_H(i)=a-d_i+R_i+x_i\ge a\). Therefore

\[x_i\ge s_i,\qquad S\ge r+2t,\qquad q_u+\rho_u\le a. \tag{5}\]

The actual degree \(x_i\) need not equal its minimum \(s_i\). No
argument below makes that substitution.

\subsection{3. The structural inequalities needed
here}\label{the-structural-inequalities-needed-here}

\textbf{Lemma 3.1.} For every selected \(ui\to w\),

\[d_i\le\rho_u+R_i,\qquad d_i\le\rho_u+\rho_w, \tag{6}\]

\[\rho_w+q_w\ge q_u-1,\qquad R_i+x_i\ge q_u+p_u. \tag{7}\]

\emph{Proof.} Every \(F\)-neighbour \(j\) of \(i\) neighbours \(u\),
because \(ui\) must dominate \(j\). At most \(\rho_u\) of these \(uj\)
edges are residual. Every other \(uj\) is selected, with a distinct
supplement \(w_j\ne w\). Since \(ui\) dominates \(w_j\) and \(u\) misses
\(w_j\), the edge \(iw_j\) exists. It is residual: \(i,w_j\) both miss
\(j\in A\), whereas every selected edge has its unique exception in
\(B\). The distinct \(w_j\) inject these incidences into residual edges
at \(i\), proving the first inequality of (6).

For the second inequality, \(uj\) must dominate \(w\), so \(jw\) exists.
Both \(j,w\) miss \(i\in A\), hence \(jw\) is residual. Distinct \(j\)
give \(d_i-\rho_u\le\rho_w\).

Every other selected label at \(u\) neighbours \(w\), because its
quasi-edge must dominate \(w\) and the supplements at \(u\) are
distinct. This proves the first inequality of (7). Finally, \(i\)
neighbours \(u\) and every missing-\(B\)-pair neighbour of \(u\) except
\(w\). There are \(q_u+p_u\) distinct vertices in that list. Hence
\(d_H(i,B)=R_i+x_i\ge q_u+p_u\). \(\square\)

In particular, every selected incidence of label \(i\) requires
\(\rho_u\ge s_i\). These proofs allow zero residual degrees and do not
invoke the residual-activity lemma. The unique-exception argument needs
no restriction on the diameter of \(H\) and no classification theorem
for its total domination number.

\subsection{4. Charging and an exact deficit
identity}\label{charging-and-an-exact-deficit-identity}

\textbf{Lemma 4.1 (inherited charging, rederived).} If \(a\ge1\), then

\[0\le s_i<a,\qquad r\ge\sum_i\frac{s_i^2}{a-s_i},\qquad t\le ca^2. \tag{8}\]

\emph{Proof.} The demand is at most \(d_i\le a-1\). For each label
choose exactly \(s_i\) of its actual selected incidences. Let \(q'_u\)
count chosen incidences at \(u\). Then \(q'_u\le a-\rho_u\), and each
chosen incidence requires \(\rho_u\ge s_i\). At a source with
\(q'_u>0\), assign charge \(\rho_u/(a-\rho_u)\) to each chosen
incidence. Its total charge is at most \(\rho_u\). Since \(x/(a-x)\) is
increasing on \([0,a)\), each of the \(s_i\) chosen incidences of label
\(i\) receives at least \(s_i/(a-s_i)\). Summing proves the residual
bound. A source with \(\rho_u=a\) has no chosen incidence and is not
divided by zero; its cost is simply unused. Zero-residual sources cause
no problem.

Using (5),

\[2t\le S-r\le\sum_i\frac{s_i(a-2s_i)}{a-s_i}.\]

For \(k=a-s>0\) the summand equals \(3a-2k-a^2/k\le(3-2\sqrt2)a\), by
\(2k+a^2/k\ge2\sqrt2a\). There are \(a\) labels, giving \(t\le ca^2\).
\(\square\)

Put \(\lambda=1-1/\sqrt2\) and \(\mathcal D=ca^2-t\). Direct algebra
yields

\[\boxed{\mathcal D=\sum_i\frac{(s_i-\lambda a)^2}{a-s_i}+\frac{S-r-2t}{2}+\frac{r-\sum_i s_i^2/(a-s_i)}2.} \tag{9}\]

All three terms are nonnegative. For clarity, the pointwise identity is

\[2ca-\frac{s(a-2s)}{a-s}=\frac{2(s-\lambda a)^2}{a-s}.\]

If \(t\ge(c-\eta)a^2\), then (9), \(a-s_i\le a\), and Cauchy--Schwarz
imply

\[\sum_i(s_i-\lambda a)^2\le\eta a^3,\quad S\le(\lambda+\sqrt\eta)a^2,\quad r\le(\lambda+\sqrt\eta-2c+2\eta)a^2. \tag{10}\]

This is stability of the \textbf{demand distribution}, not a graph
edit-distance or near-bipartiteness theorem.

\Needspace{14\baselineskip}
\subsection{5. Shared pair capacity and a quadratic
loss}\label{shared-pair-capacity-and-a-quadratic-loss}

\textbf{Lemma 5.1 (good-label pair capacity).} Let \(I\subseteq A\) have
\(s_i\ge La\) for \(i\in I\), where \(L>0\), and assume
\(|A\setminus I|\le\theta a\). Set

\[Z=\{u\in B:\rho_u\ge La\},\quad z=|Z|,\quad T=(L+\theta)a+1,\quad P=\sum_{i\in I}x_i.\]

If \(k\) sources in \(Z\) have \(q_u>T\), then

\[P\le(z-k)T+kz-\frac{k(k+1)}2\le\frac{z^2+T^2}{2}. \tag{11}\]

\emph{Proof.} All selected incidences with labels in \(I\) have sources
in \(Z\). Outside \(Z\), a vertex has \(\rho_u<La\) and can select only
labels outside \(I\), so \(q_u\le\theta a\) and
\(d_H(u,A)<(L+\theta)a\). A source with \(q_u>T\) therefore cannot use a
supplement outside \(Z\), by (7).

Let \(U\subseteq Z\) be the \(k\) busy sources. The other \(z-k\)
sources supply at most \((z-k)T\) incidences in \(I\). All outgoing arcs
of \(U\) lie inside \(Z\) and occupy distinct unordered pairs incident
with \(U\). Their number is at most

\[\binom z2-\binom{z-k}2=kz-\frac{k(k+1)}2.\]

This bounds all their outgoing arcs, and hence their incidences in
\(I\). We have used one common pair budget, not a separate favourable
capacity for each source. Finally,

\[(z-k)T+kz-\frac{k(k+1)}2=\frac{z^2+T^2}{2}-\frac{(z-T-k)^2}{2}-\frac k2,\]

proving (11). It holds even when \(T\) is not integral or exceeds \(z\).
\(\square\)

\textbf{Proof of Theorem A.} Suppose \(a\ge25\) and \(t\ge c_1a^2\). Use

\[\eta=\frac1{5000},\quad L=\frac6{25},\quad \theta=\frac9{125}.\]

Since \(\lambda-L>0\) and \((\lambda-L)^2>1/360\), (10) implies that
fewer than \(\theta a\) labels have \(s_i<La\). Take \(I\) to be the
remaining labels. Thus

\[P\ge La|I|\ge\frac{696}{3125}a^2. \tag{12}\]

Here \(\sqrt\eta=\sqrt2/100\). The residual bound (10) becomes

\[r\le\left(-\frac{4999}{2500}+\frac{151}{100}\sqrt2\right)a^2<\frac{17}{125}a^2.\]

Each source in \(Z\) costs at least \(La\), giving \(z\le17a/30\). Also
\(a\ge25\) implies

\[T=(L+\theta)a+1\le\frac{44}{125}a.\]

Apply (11):

\[P\le\frac12\left[\left(\frac{17}{30}\right)^2+\left(\frac{44}{125}\right)^2\right]a^2=\frac{250321}{1125000}a^2<\frac{696}{3125}a^2,\]

contradicting (12). The exact gap is \(239a^2/1125000\). This proves
(1). \(\square\)

All irrational comparisons above are elementary rational-square
comparisons. Two separately written checkers verify the constants, one
in \(\mathbb Q(\sqrt2)\) and one using rational enclosures for
\(\sqrt2\); neither numerical optimisation nor finite graph enumeration
is needed by the proof.

\subsection{6. The all-order maximum-degree
consequence}\label{the-all-order-maximum-degree-consequence}

\textbf{Proof of Corollary A.} A universal vertex forces an
edge-critical graph to be a star: an edge between other vertices would
be redundant. For \(n\ge4\) that star has fewer than
\(\lfloor n^2/4\rfloor\) edges. Hence take \(a\ge1\).

Suppose \(b\ge\beta_1n\) and \(m\ge\lfloor n^2/4\rfloor\). Put
\(\varepsilon=0\) for even \(n\) and \(1\) for odd \(n\). Then

\[t\ge(b-n/2)^2-\varepsilon/4.\]

The function \(((x-1/2)/(1-x))^2\) is increasing on \([1/2,1)\) and
equals \(c_1\) at \(x=\beta_1\). Consequently

\[t\ge c_1(n-b)^2-\varepsilon/4=c_1a^2+c_1(2a+1)-\varepsilon/4>c_1a^2,\]

where \(c_1>1/12\) gives the final strictness. This contradicts Theorem
A for \(a\ge25\).

For \(1\le a\le24\), the exact lower bound \(\beta_1>613168/10^6\)
implies

\[b\ge b_0(a)=\left\lceil\frac{613168(a+1)}{386832}\right\rceil>a+1.\]

Since \(n=a+b+1\),

\[t\ge\left\lfloor\frac{(b-a-1)^2}{4}\right\rfloor\ge t_0(a)=\left\lfloor\frac{(b_0(a)-a-1)^2}{4}\right\rfloor.\]

The following 24 hand-checkable cases have \(t_0(a)>ca^2\),
contradicting (8):

\Needspace{16\baselineskip}
\begin{longtable}[]{@{}rrrrrr@{}}
\toprule\noalign{}
\(a\) & \(b_0\) & \(t_0\) & \(a\) & \(b_0\) & \(t_0\) \\
\midrule\noalign{}
\endhead
\bottomrule\noalign{}
\endlastfoot
1 & 4 & 1 & 13 & 23 & 20 \\
2 & 5 & 1 & 14 & 24 & 20 \\
3 & 7 & 2 & 15 & 26 & 25 \\
4 & 8 & 2 & 16 & 27 & 25 \\
5 & 10 & 4 & 17 & 29 & 30 \\
6 & 12 & 6 & 18 & 31 & 36 \\
7 & 13 & 6 & 19 & 32 & 36 \\
8 & 15 & 9 & 20 & 34 & 42 \\
9 & 16 & 9 & 21 & 35 & 42 \\
10 & 18 & 12 & 22 & 37 & 49 \\
11 & 20 & 16 & 23 & 39 & 56 \\
12 & 21 & 16 & 24 & 40 & 56 \\
\end{longtable}

For each fixed \(a\) the lower bound is increasing in \(b\ge b_0(a)\),
so this covers all orders, not just the 24 minimum orders in the table.
A rational upper enclosure \(c<3/2-1414213562373095/10^{15}\) suffices
to check every entry. \(\square\)

\textbf{Audit of the older coefficient.} The same argument with \(c\)
and \(\beta_0\) proves the inherited theorem directly for all \(a\ge1\),
because \(3c>1/4\). This confirms its floor/odd-order strictness without
the residual-activity lemma. No blocking error was found in that
rederivation. It is still a candidate pending independent review.

\textbf{A scalar comparison only.} At \(a=7000,b=11200,h=2000,z=2800\),
take \(z\) residual rows of degree \(h\), the other rows of degree one,
and every label demand \(h\). Then \(r=5608400\), \(S=14000000\),
\(t=(S-r)/2=4195800\). These numbers satisfy both basic charging and the
active scalar charging equality \(r-b=\sum s_i(s_i-1)/(a-s_i)\), and
have \(c_1a^2<t<ca^2\). This witnesses a genuine strengthening of the
scalar inequality. It is not a graph, and is not asserted to pass the
older source--supplement tests.

\subsection{7. A general spare-source
theorem}\label{a-general-spare-source-theorem}

Let \(h\ge1\) be the maximum residual degree and \(Z=\{u:\rho_u=h\}\)
with \(|Z|=z\). Let \(I=\{i:s_i=h\}\) have size \(\ell\). If \(z<h\)
then \(I\) is empty by (6), so assume \(h\le z\le h^2+h\).

Every selected incidence of a label in \(I\) has its source in \(Z\),
and \(h\le x_i\le z\). At such an incidence (6) gives \(d_i\le2h\).
Since \(R_i=d_i-h\), call \(i\) \emph{exceptional} if \(d_i=2h\), and
\emph{ordinary} otherwise. An exceptional selected incidence must have
its supplement in \(Z\), by the second inequality of (6). That
supplement misses the label and cannot be one of its selected sources.
Thus \(x_i\le z-1\) for an exceptional label.

Writing \(D_i=R_i+x_i\), we have

\[D_i\le h+x_i-1\text{ for ordinary }i,\qquad D_i=h+x_i\text{ for exceptional }i,\]

and in either case \(D_i\le h+z-1\).

For \(u\in Z\), let \(y_u\) count outgoing incidences in \(I\), \(e_u\)
count the exceptional ones, and \(p^E_u\) count incoming exceptional
arcs. Let \(\sigma_u=q_u+p_u\) be its actual pair degree. Define the
weighted label load

\[L_u=\sum_{i\in I:\,ui\text{ selected}}\frac h{x_i}.\]

Each label contributes \(h\) in total. Also

\[y_u\le q_u\le\sigma_u-p^E_u,\quad p^E_u\le z-1,\quad \sum_{u\in Z}e_u=\sum_{u\in Z}p^E_u. \tag{13}\]

The last equality holds because every exceptional arc is wholly inside
\(Z\).

\Needspace{5\baselineskip}
\textbf{Weighted local bound.} We claim

\[L_u\le2h-1+\frac{e_u-p^E_u}{h+1}. \tag{14}\]

If \(y_u=0\), then \(L_u=e_u=0\), while \(p^E_u\le z-1\le h^2+h-1\)
makes the right side nonnegative. If \(y_u>0\) and \(\sigma_u\le2h-1\),
each weight is at most one, so

\[L_u\le y_u\le2h-1-p^E_u\le2h-1+\frac{e_u-p^E_u}{h+1}.\]

It remains to take \(\sigma_u=2h+j\). The endpoint inequality (7) bounds
the actual degree by \(h+z-1\), hence \(0\le j\le h^2-1\). Every
ordinary label selected at \(u\) has \(x_i\ge h+j+1\), and every
exceptional label has \(x_i\ge h+j\). Thus

\[L_u\le\frac{h(2h+j-p^E_u)}{h+j+1}+\frac{he_u}{(h+j)(h+j+1)}. \tag{15}\]

Subtract (15) from the right side of (14). The result is

\[\frac{(h-1)(j+1)}{h+j+1}+p^E_u\left(\frac h{h+j+1}-\frac1{h+1}\right)+e_u\left(\frac1{h+1}-\frac h{(h+j)(h+j+1)}\right).\]

Every term is nonnegative: the middle coefficient is
\((h^2-j-1)/[(h+j+1)(h+1)]\), and the final numerator after taking a
common denominator is \(j(2h+j+1)\). This proves (14) for every source,
including those with no selected label in \(I\).

Sum (14) over \(Z\). The exceptional incoming and outgoing charges
cancel by (13), leaving \(h\ell\le(2h-1)z\). This proves Theorem B.
\(\square\)

\textbf{Corollaries.} For \(z=h+k\) with \(0\le k\le h^2\),

\[\ell\le\left\lfloor\frac{(2h-1)(h+k)}h\right\rfloor.\]

For \(1\le k\le h\) this is \(\ell\le2h+2k-2\). In particular \(k=1\)
gives the older one-spare-source bound, and \(k=2\) gives \(2h+2\) for
\(h\ge2\). The case \(h=1,k=2\) lies outside this theorem; no conclusion
for it is silently asserted. For \(k=0\) the bound is \(2h-1\).

\textbf{Sharpness of reduced premises.} Set
\(\ell=\lfloor(2h-1)z/h\rfloor\). Number sources modulo \(z\) and give
label \(i\) the distinct sources \(ih,ih+1,\ldots,ih+h-1\) modulo \(z\).
All labels are ordinary with \(x_i=h,R_i=h-1,d_i=2h-1\). Each source
appears at most \(2h-1\) times, so the endpoint inequalities hold with
no exceptional arcs. These data attain the integer label bound for the
reduced ordinary incidence model. They do not specify the remaining
residual graph, missing-pair structure, or graph \(G\), so this is not
an extremal-graph construction.

The cancellation is essential. An exceptional source can locally exceed
\(2h-1\) in weighted load; suppressing the compensation term would be an
invalid proof. The abstract tests include such systems rather than
testing only the ordinary case.

\subsection{8. Validation, formalisation and remaining
obligations}\label{validation-formalisation-and-remaining-obligations}

The code and data in this checkpoint make the checks repeatable.
\texttt{src/exact\_checks.py} verifies the quadratic-field constants, 24
small-\(a\) cases, 25,755 pair identities, 5,050 deficit identities, and
338,350 spare-coefficient triples. \texttt{src/interval\_check.py}
independently verifies the constants and rounding with rational
enclosures and rejects three corrupted parameter choices. Neither is a
formal proof for all graph orders.

Actual-graph regression reconstructs 4,694 selected systems from 781
distinct labelled critical graphs. It includes all 21 critical graphs
found in NetworkX's atlas, 120 generated larger critical graphs, four
cycle expansions, and inherited examples. The atlas criticality
decisions were cross-checked using a separate distance-based
implementation. Some roots' quasi-edge choices were sampled; details and
seeds are saved. There are 191 nonempty-demand systems and 194 nonempty
good-label pair tests, but \textbf{no positive-surplus example and no
nonempty maximum-demand spare class in the theorem's range}. Those
limitations prevent claiming a nonvacuous actual-graph test of the new
dense contradiction or spare theorem.

To exercise the latter's premises directly, the abstract tests accept
7,435 reduced incidence systems, of which 3,555 have exceptional arcs.
They also exhaust 59,809 oriented simple graphs through five vertices,
making 358,058 pair-capacity threshold checks. These are not critical
graphs; they test the counting mechanisms only. All recorded tests pass.
Omitted, duplicate and self-supplement graph selections are rejected.

\textbf{Formal scope.} \texttt{formal/QuasiCore.lean} has been checked
using official Lean 4.19.0 in GitHub run 34166620020, at source commit
\texttt{93a7cdc5ff2a3677d9985941a5df0c4ad08bd1f6}. The five checked
propositions concern uniqueness of an exception, forcing a neighbour
when the source misses a vertex, and the common-miss obstruction to
selecting a forced cross-edge. The file assumes explicit quasi-edge
predicates; it does not prove they exist for every critical graph. It
contains no \texttt{sorry}, new axiom declaration or
\texttt{native\_decide}. The printed axiom dependencies are either empty
or the standard \texttt{propext}, \texttt{Classical.choice},
\texttt{Quot.sound}. \textbf{No claim of a formally verified charging
theorem, spare-source theorem, maximum-degree bound or n=28 proof
follows.}

The general formalisation literature already includes an end-to-end
Lean/SAT graph-generation framework {[}3{]}. This small local slice is
not claimed as the first formal treatment of Murty--Simon. The next
useful formal targets are quasi-edge existence and the finite injections
behind (6), then the charging sum.

The important independent-review questions are: whether the
selected-edge injections are genuinely injective and residual; whether
all actual selected incidences of good labels are confined to \(Z\);
whether the shared unordered-pair count in (11) covers every
orientation; and whether exceptional arcs in Theorem B have both
endpoints in \(Z\), enabling exact cancellation. Paul is undertaking
specialist outreach. No external endorsement, whole-order extension or
resolution of the general conjecture is claimed.

\subsection{References}\label{references}

{[}1{]} T. W. Haynes, M. A. Henning, L. C. van der Merwe and A. Yeo,
\emph{A maximum degree theorem for diameter-2-critical graphs}, Central
European Journal of Mathematics 12 (2014), 1882--1889. DOI:
10.2478/s11533-014-0449-3. Author-institution record:
\url{https://dc.etsu.edu/etsu-works/15862/}.

{[}2{]} A. Jabalameli, A. Behjati, M. Saghafian, M. M. Shokri, M.
Ferdosi and S. Bahariyan, \emph{Improving the Bounds On Murty\_Simon
Conjecture}, arXiv:1610.00360, canonical v2 dated 10 October 2016.
\url{https://arxiv.org/abs/1610.00360}. Related author-posted full text
is identified and distinguished in \texttt{literature/COMPARISON.md}.

{[}3{]} M. Kirchweger, P. Manrique and S. Szeider, \emph{Formally
Verified Graph Generation with SAT Modulo Symmetries and Lean}, IJCAR
2026, LNCS 16688, 117--135. Published 24 July 2026. DOI:
10.1007/978-3-032-32589-1\_8.
\url{https://doi.org/10.1007/978-3-032-32589-1_8}.

{[}4{]} Paul Lenz project, ChatGPT development, \emph{General-order
residual charging and coupled source--supplement cuts}, 7 September
2026, \texttt{2026-09-07-coupled-resource-v2/PROOF.md}, Git blob
\texttt{3d3b4ab5b2da8207f17a3a2afcce49e7119d78e2}. Candidate project
source, not an independently published theorem.

{[}5{]} Paul Lenz project, ChatGPT development, \emph{LocalIncidence
v7}, 7 September 2026, original
\texttt{MurtySimon\_GeneralN\_LocalIncidence\_v7.zip}, Section 7; its
archive hash is recorded in \texttt{PROVENANCE.json}. Candidate
one-spare-source antecedent.

\end{document}
